\begin{figure}[htbp]
\centering
\begin{tikzpicture}[font=\small, x=0.95cm, y=0.55cm, >=Latex]
\fill[blue!9] (0,2) rectangle (5,4);
\fill[teal!9] (0,0) rectangle (5,2);
\fill[blue!30] (0,2) rectangle (5,3);
\fill[teal!30] (0,1) rectangle (5,2);
\draw[step=1, gray!60] (0,0) grid (5,4);
\draw[thick] (0,2) -- (5,2);
\node[anchor=east] at (-0.2,3) {Rank $r$};
\node[anchor=east] at (-0.2,1) {Rank $r+1$};
\draw[->,semithick,blue!70!black] (5,2.5) -- (5.5,2.5) -- (5.5,0.55) -- (5.85,0.55);
\draw[->,semithick,teal!70!black] (5,1.5) -- (5.25,1.5) -- (5.25,3.45) -- (5.85,3.45);
\node[anchor=west,align=left] at (5.85,3.45) {Halo de $r$};
\node[anchor=west,align=left] at (5.85,0.55) {Halo de $r+1$};
\end{tikzpicture}
\caption{Duas faixas consecutivas. Cada rank recebe a linha de fronteira do vizinho, incluindo os estados necessários às consultas diagonais.}
\label{fig:tb2-faixas}
\end{figure}
